\documentclass[11pt,a4paper]{article}
\usepackage[utf8]{inputenc}
\usepackage[T1]{fontenc}
\usepackage{amsmath}
\usepackage[a4paper,margin=2.5cm]{geometry}
\usepackage{xcolor}
\usepackage{listings}
\usepackage{booktabs}
\usepackage{enumitem}
\usepackage{longtable}
\usepackage[hidelinks]{hyperref}

\definecolor{codebg}{rgb}{0.96,0.96,0.96}
\definecolor{kw}{rgb}{0.0,0.3,0.6}
\definecolor{cmt}{rgb}{0.3,0.5,0.3}
\lstset{
  backgroundcolor=\color{codebg},
  basicstyle=\ttfamily\small,
  keywordstyle=\color{kw}\bfseries,
  commentstyle=\color{cmt}\itshape,
  breaklines=true,
  frame=single,
  framesep=4pt,
  showstringspaces=false,
  columns=fullflexible,
}

\title{\textbf{Praktikum: Secret-Key-Generation (SKG) Sederhana\\ Alice (Laptop) $\leftrightarrow$ Bob (ESP32)}}
\author{Modul Praktikum}
\date{}

\begin{document}
\maketitle

\section{Tujuan}
Mahasiswa mampu:
\begin{enumerate}[noitemsep]
  \item Memahami prinsip \emph{physical-layer secret key generation} dari sinyal RSSI.
  \item Menjalankan protokol SKG antara laptop (Alice) dan ESP32 (Bob).
  \item Membuktikan kedua perangkat menghasilkan \textbf{kunci yang sama} tanpa mengirim kunci.
\end{enumerate}

\section{Konsep Singkat}
Dua perangkat mengukur RSSI pada kanal WiFi yang sama. Karena resiprositas kanal, hasil
pengukuran berkorelasi. Dari korelasi ini kunci diekstraksi melalui empat langkah:

\begin{longtable}{p{3.2cm} p{10.5cm}}
\toprule
\textbf{Langkah} & \textbf{Penjelasan} \\
\midrule
1. Quantisasi guard-band & RSSI diubah menjadi bit. Sampel ambigu di pita $\mu \pm \alpha\sigma$ dibuang. \\
2. Sifting & Tukar \emph{posisi} sampel yang dipakai, ambil irisan. Posisi aman dibagi publik. \\
3. Rekonsiliasi (Cascade) & Perbaiki bit yang berbeda memakai \emph{parity} + \emph{binary search} sampai bit identik. \\
4. Privacy amplification & Kunci akhir $=$ SHA3-256(bit). Menghapus informasi parity yang sempat bocor. \\
\bottomrule
\end{longtable}

\noindent Verifikasi akhir: kedua sisi menukar $H(H(\text{kunci}))$ untuk mengecek \textbf{MATCH}.

\paragraph{Keamanan.} Yang melewati kanal publik hanya posisi indeks dan parity, bukan nilai
RSSI. Penyadap (Eve) tidak memiliki nilai RSSI, dan privacy amplification menghapus parity
yang bocor, sehingga Eve tidak dapat menurunkan kunci.

\section{Alat dan Bahan}
\begin{itemize}[noitemsep]
  \item \textbf{Laptop (Alice):} Python 3.8+ (tanpa pustaka tambahan).
  \item \textbf{ESP32 (Bob):} Arduino IDE atau \texttt{arduino-cli}, board package \textbf{esp32}, kabel USB.
  \item \textbf{Jaringan:} satu hotspot WiFi \textbf{2.4 GHz} untuk laptop dan ESP32
        (disarankan hotspot HP; hindari WiFi kampus karena \emph{client isolation}).
\end{itemize}

\section{Struktur Berkas}
\begin{lstlisting}
mahasiswa-skg/
  alice/   (laptop)  alice.py, skg_common.py, sha3_256.py, pcg32.py, net.py, synced_alice.csv
  bob/     (ESP32)   bob_esp32/bob_esp32.ino, skg.h, synced_bob.h
  docs/              praktikum.tex
\end{lstlisting}

\section{Langkah Praktikum}

\subsection{A. Siapkan Jaringan}
\begin{enumerate}[noitemsep]
  \item Nyalakan hotspot HP (2.4 GHz). Catat SSID dan password.
  \item Sambungkan laptop ke hotspot tersebut.
\end{enumerate}

\subsection{B. Bob (ESP32)}
\begin{enumerate}[noitemsep]
  \item Tambahkan URL board (Arduino IDE $\to$ Preferences):
  \begin{lstlisting}
https://espressif.github.io/arduino-esp32/package_esp32_index.json
  \end{lstlisting}
  \item Boards Manager $\to$ install \textbf{esp32}.
  \item Buka \texttt{bob/bob\_esp32/bob\_esp32.ino}, isi WiFi:
  \begin{lstlisting}[language=C++]
const char* WIFI_SSID = "NAMA_HOTSPOT";
const char* WIFI_PASS = "PASSWORD_HOTSPOT";
const double ALPHA    = 1.0;   // harus sama dengan alice.py
  \end{lstlisting}
  \item Pilih Board \textbf{ESP32 Dev Module} dan Port, klik \textbf{Upload}.
  \item Buka Serial Monitor (115200), catat \texttt{IP ESP32}.
\end{enumerate}

\noindent Alternatif \texttt{arduino-cli}:
\begin{lstlisting}
arduino-cli core install esp32:esp32
arduino-cli compile --fqbn esp32:esp32:esp32 bob/bob_esp32
arduino-cli upload  --fqbn esp32:esp32:esp32 -p COM_ESP32 bob/bob_esp32
arduino-cli monitor -p COM_ESP32 -c baudrate=115200
\end{lstlisting}

\subsection{C. Alice (Laptop)}
Edit \texttt{alice/alice.py}:
\begin{lstlisting}[language=Python]
BOB_IP = "192.168.x.y"   # IP ESP32 dari Serial Monitor
ALPHA  = 1.0             # harus sama dengan bob_esp32.ino
\end{lstlisting}

\subsection{D. Jalankan}
Pastikan ESP32 sudah ``menunggu Alice''. Lalu di laptop:
\begin{lstlisting}
cd alice
python alice.py
\end{lstlisting}

\subsection{E. Hasil yang Benar}
\begin{lstlisting}
# Laptop (Alice)
Sifting: posisi bersama = 123
Kunci Alice : 4c032a3015f75237...067ed8
Status      : MATCH

# Serial Monitor (Bob / ESP32)
Rekonsiliasi selesai, query parity = 88
STATUS: MATCH
Kunci Bob: 4c032a3015f75237...067ed8
\end{lstlisting}
\textbf{Kunci Alice $=$ Kunci Bob $\Rightarrow$ praktikum berhasil.}

\section{Troubleshooting}
\begin{longtable}{p{4.8cm} p{8.6cm}}
\toprule
\textbf{Gejala} & \textbf{Solusi} \\
\midrule
ESP32 cetak titik terus & WiFi salah / 5 GHz. Cek SSID/pass, pakai hotspot 2.4 GHz. \\
Alice \texttt{Gagal connect} & IP ESP32 salah atau beda WiFi. Samakan jaringan, perbarui \texttt{BOB\_IP}. \\
Alice menggantung lalu gagal & \emph{client isolation}. Gunakan hotspot HP. \\
\texttt{STATUS: RETRY} & Kanal bising. Naikkan \texttt{ALPHA} (mis. 1.2) di kedua sisi. \\
Port COM tak terbuka & Tutup Serial Monitor sebelum upload. \\
\bottomrule
\end{longtable}

\section{Pertanyaan Diskusi}
\begin{enumerate}[noitemsep]
  \item Mengapa menukar \emph{posisi} indeks aman, tetapi menukar nilai bit tidak?
  \item Apa pengaruh menaikkan $\alpha$ terhadap jumlah bit dan BER?
  \item Mengapa privacy amplification (SHA3) diperlukan setelah rekonsiliasi?
  \item Hitung entropi bersih bila bit bersama $=123$ dan parity bocor $=88$.
\end{enumerate}

\end{document}
